% =============================================================================
% GroundLM 2026 Shared Tasks - GoldenViewVQA system description paper
% Official ACL style files, copied UNMODIFIED. Do not alter margins, spacing,
% fonts or page dimensions.
% Compile: latexmk -pdf main.tex
% =============================================================================
\documentclass[11pt]{article}
\usepackage[final]{acl}

\usepackage{times}
\usepackage{latexsym}
\usepackage[T1]{fontenc}
\usepackage[utf8]{inputenc}
\usepackage{microtype}
\usepackage{graphicx}
\usepackage{booktabs}
\usepackage{multirow}
\usepackage{amsmath}
\usepackage{url}

% Set the team name ONCE. Must be byte-identical to the evaluator Space
% registration and the OpenReview `teamname` field. Immutable once registered.
\newcommand{\teamname}{Ensar}

\title{\teamname{} at GroundLM 2026 Shared Tasks: \\
       Symbolic Scene Serialisation for View-Level Evidence Identification}

\author{
  Mostafa Gaafar\thanks{\ Equal contribution.} \\
  The American University \\ in Cairo, Egypt \\
  \texttt{mostafa21314@aucegypt.edu} \\
  \And
  Omar Saqr\footnotemark[1] \\
  The American University \\ in Cairo, Egypt \\
  \texttt{Omar\_saqr@aucegypt.edu} \\
}

\begin{document}
\maketitle

\begin{abstract}
GoldenViewVQA asks a system to name which of six synchronised vehicle cameras
carries the evidence for a question, and to answer that question. We submit a
two-stage system built on a symbolic scene serialisation: a nuScenes-fine-tuned
detector is run over all six views, detections are back-projected onto the
ground plane, associated across cameras into objects carrying an explicit
view-set, and rendered as text. A language model then predicts the supporting
view and the answer from that text alone (System~1), or from the text together
with the six raw images (System~2). On the development split, System~1 reaches
$0.3816$ view macro-accuracy and System~2 reaches $0.6786$, against $0.1429$
for an always-\textsc{cam\_front} baseline. A paired bootstrap shows the
pixels contribute a significant gain to view attribution
($+0.297$ macro, 95\% CI $[+0.078,+0.461]$) and \emph{no} gain to answering
($-0.018$, CI $[-0.055,+0.000]$). Our most consequential finding is negative:
answer accuracy is already $0.9455$ with no visual input at all, so the answer
head is largely solvable from option phrasing and essentially all measurable
signal in this benchmark lies in the view head.
\end{abstract}

% =============================================================================
\section{Introduction}
\label{sec:intro}

GoldenViewVQA \citep{wang2026doesanswercomefrom} evaluates whether a
multi-view model can say \emph{where} its evidence came from. Given six
synchronised nuScenes \citep{caesar2020nuscenes} camera images and a
four-option multiple-choice question, a system must predict the supporting
camera view --- one of the six channels or \textsc{none\_of\_the\_above} ---
and the answer. Submissions are ranked by joint accuracy, which requires both.

Two properties of the released data shape any reasonable approach. First,
there is no training split: 55 labelled development records and 59 input-only
test records, and nothing else. This is an inference-time task, not a training
task. Second, the view labels are severely imbalanced. In development,
\textsc{cam\_front} accounts for $38/55$ records while two classes hold a
single example each. Because view macro-accuracy averages recall over the
seven classes present in gold, each class carries roughly $14.3$ points, and a
single flipped prediction on a singleton class moves the headline metric by
that much.

Our approach is deliberately not an end-to-end model. Serialising an
object-level scene representation for a language model, rather than feeding it
pixels, follows \citet{chen2024driving}, who fuse an object-level vector
modality with an LLM for explainable driving, and \citet{mao2023gptdriver},
who cast driving decisions as language modelling over a textual scene. Driving
question answering has likewise been framed as structured reasoning over scene
graphs \citep{sima2024drivelm}. What distinguishes GoldenViewVQA from these is
that the label is the \emph{evidence source} rather than the decision, which
makes per-object camera provenance --- not just object presence --- the
quantity a serialisation must preserve.

We therefore detect agents across all six views and associate them in
bird's-eye view following \citet{gaafar2025geodrive}, and serialise the result
so that each object carries the set of cameras it was seen in. The reasoning is
that if every detected object is attributed to the specific cameras it appears
in, then view selection reduces to grounding the question to an object and
reading off its view-set. We test that reduction directly, and then test what
raw pixels add on top of it.

Our contributions are:

\begin{itemize}
  \item A symbolic serialisation of a surround-view scene that carries
        explicit per-object view-sets, cross-camera identity, and negative
        evidence, together with two systems built on it.
  \item Measured evidence that the pixels help view attribution significantly
        and do not help answering at all.
  \item A negative result that we believe matters for the benchmark: answer
        accuracy reaches $0.9455$ with no images, which implies the option
        text leaks the answer and that joint accuracy is, in practice, a
        view-accuracy metric.
  \item An error analysis of the perception-to-text interface, including a
        failure mode where a missed detection is serialised as confirmed
        negative evidence.
\end{itemize}

% =============================================================================
\section{Data and Preprocessing}
\label{sec:data}

\paragraph{Released data.} The task provides \texttt{eval.jsonl} (55 labelled
records), \texttt{eval\_inputs.jsonl} and \texttt{test\_inputs.jsonl} (59
records, labels withheld), machine-readable input and submission schemas, and
an official evaluator. Each record carries a question, four options, a
\texttt{question\_group} in \{\textit{causality}, \textit{counterfactual},
\textit{intent\_prediction}\}, a nuScenes \texttt{sample\_token}, and a
\texttt{views} map from camera name to a nuScenes-relative image path.
Development is causality-heavy ($27/55$) while test is more evenly spread and
leans more on intent prediction ($20/59$), so the two splits are not
distributionally interchangeable.

\paragraph{Images.} Images are not redistributed. We obtained nuScenes
\texttt{v1.0-trainval} camera data under the nuScenes Terms of Use. All 114
records resolve to 438 unique frames across 66 scenes; we verified $438/438$
resolve locally before running anything. Only keyframe (\texttt{samples/})
camera data is required --- no sweeps, lidar or radar --- which reduces the
download substantially.

\paragraph{No ground-truth annotations at inference.} The task permits
inferring from the images but not reading nuScenes annotations. We therefore
removed every annotation-derived component from the GeoDrive pipeline:
3D bounding boxes, ego pose, the OpenStreetMap and GPS map lookups, and
ground-truth future trajectories. The only nuScenes-derived quantity we retain
is the camera rig geometry, which we treat as published hardware
specification: across all 850 trainval logs, focal length varies by
$\leq\!1.3\%$ and mount positions by under $6$\,cm, so we bake per-camera
intrinsics and extrinsics into constants. Inference reads six JPEG files and
touches no nuScenes table.

\paragraph{No augmentation.} No data augmentation is used. There is no
training stage in which it could act.

% =============================================================================
\section{System Architecture}
\label{sec:system}

\subsection{Shared perception front-end}

The front-end converts six images into a list of objects, each annotated with
the set of cameras it was detected in.

\paragraph{Detection.} We run YOLOv8-m \citep{ultralytics2023yolov8},
fine-tuned on nuScenes over seven agent classes (car, truck, bus, trailer,
motorcycle, bicycle, pedestrian) as in GeoDrive
\citep{gaafar2025geodrive}, independently over each of the six images. We set
the confidence threshold to $0.15$, deliberately low: a missed object is
unrecoverable in the serialised text, whereas a false positive can still be
discounted downstream by a model that sees pixels.

\paragraph{Ground-plane projection.} For each detection we back-project the
bottom-centre of the box through the fixed rig geometry onto the plane
$z=0$ in the ego frame, giving a metric position. Rays that do not meet the
plane in front of the camera are discarded.

\paragraph{Cross-camera association.} Detections of the same class within
$2.0$\,m in bird's-eye view --- one vehicle width, following
\citet{gaafar2025geodrive} --- are merged into a single object. Critically,
and unlike the GeoDrive implementation, a merged duplicate is \emph{absorbed}
rather than discarded, so the camera it came from is recorded. The resulting
\texttt{views} list is the quantity the task actually asks about; discarding it
makes the full rendering below impossible.

\subsection{Serialisation}

Objects are rendered as text in two forms. The \textbf{full render} emits one
section per camera in fixed order, objects sorted near-to-far, with one shared
identifier per object across every section it appears in, explicit
cross-references (\texttt{[same object also in CAM\_FRONT]}) generated from the
view-set, and an explicit \texttt{no agents detected} line for empty cameras.
The \textbf{isolated render} emits a single camera with cross-references
stripped and identifiers local to that view.

Bearing is reported in degrees, since azimuth is the well-conditioned axis of
monocular geometry and is what view attribution depends on. Range is quantised
into bands (near $<\!10$\,m, mid $10$--$25$\,m, far $>\!25$\,m) rather than
reported as a point estimate, because flat-ground monocular range error grows
quadratically with distance.

\subsection{Submitted systems}

\paragraph{System 1 (text only).} The full render, the question and the four
options are given to a language model, which returns the predicted view and
answer. No images are passed. This isolates what the symbolic layer
contributes.

\paragraph{System 2 (text and pixels).} The full render is presented first,
followed by all six images, each labelled with its camera name. The prompt
states that the detector cannot see traffic lights, signs, cones, road
markings, lane geometry or barriers, and instructs the model to trust the
pixels where text and pixels disagree.

\paragraph{Released prompts.} Prompting has two layers, and we release both.
The per-record prompt is generated by our sample builders from the serialised
scene, the question and the options. The runner instructions --- how records are
worked through, how text is weighed against pixels, and the output contract ---
are released verbatim in the \texttt{prompts/} directory of our code, one file
per system.

\paragraph{Models and compute.} We ran the serialisation through two hosted
reasoning models: Claude Opus~5 with extended thinking, and GPT-5.6 Sol at its
highest reasoning setting. Within each model the prompt is held fixed across
systems, so the System~2 versus System~1 contrast measures the contribution of
pixels rather than a change of model. Our best official submission used Claude
Opus~5. No model is trained or fine-tuned by us; the only learned weights are
the inherited detector checkpoint. Detection ran on a single consumer GPU in
under two minutes for all 438 images. The 114 records were processed in shards
of 9--11 records to bound context per call.

\paragraph{System 3 (per-view evidence ranking).} We additionally implemented an
approximation of an information-gain formulation. Each camera is presented on
its own, with the isolated render and that camera's image, and the model reports
an answer and an ordinal evidence-strength score for that view; a question-only
pass supplies a prior answer. The predicted view is the argmax of evidence
strength. This is an approximation and we label it as such: the hosted
interfaces expose no token log-probabilities, so we cannot compute the
distributional divergence the formulation calls for, and an ordinal judgment
elicited in natural language is a coarse substitute. It costs seven passes per
question instead of one.

% =============================================================================
\section{Experimental Setup}
\label{sec:setup}

We report the four official metrics: \texttt{view\_accuracy} (micro),
\texttt{view\_macro\_accuracy} (mean per-view recall over classes present in
gold), \texttt{answer\_accuracy}, and \texttt{joint\_accuracy}. All numbers
come from the organizers' \texttt{evaluate.py}, imported rather than
reimplemented, so our reported metrics cannot drift from the official
definition.

\paragraph{Baselines.} We include four label-only baselines that read neither
text nor pixels: \textit{constant} (always \textsc{cam\_front}, answer A),
\textit{majority} (most frequent gold view and answer), \textit{prior\_random}
(sampled from the development label distribution), and
\textit{uniform\_random}.

\paragraph{Uncertainty.} With $n=55$, point estimates are close to
meaningless, so we report 95\% percentile bootstrap intervals over questions
($2000$ resamples). For system comparisons we use a \emph{paired} bootstrap
that resamples the same questions for both systems, so per-question difficulty
cancels. To calibrate how much of an observed difference is noise, we ran
\textit{uniform\_random} across 20 seeds, changing nothing else: macro
accuracy had mean $0.112$, standard deviation $0.069$, and range
$[0.015,0.223]$. A $0.208$-wide spread arises from the random seed alone,
which is wider than the gap between several of our baselines.

% =============================================================================
\section{Results}
\label{sec:results}

\begin{table}[t]
\centering\small
\begin{tabular}{lcccc}
\toprule
System & View & Macro & Ans. & Joint \\
\midrule
\multicolumn{5}{l}{\textit{label-only baselines}} \\
uniform\_random   & 0.146 & 0.223 & 0.327 & 0.091 \\
prior\_random     & 0.491 & 0.102 & 0.309 & 0.200 \\
constant          & 0.691 & 0.143 & 0.255 & 0.182 \\
majority          & 0.691 & 0.143 & 0.309 & 0.200 \\
\midrule
\multicolumn{5}{l}{\textit{Claude Opus 5}} \\
System 1 (text)    & 0.709 & 0.382 & 0.945 & 0.709 \\
System 2 (+pixels) & 0.873 & 0.679 & 0.927 & 0.855 \\
\midrule
\multicolumn{5}{l}{\textit{GPT-5.6 Sol}} \\
System 1 (text)    & 0.727 & 0.481 & 0.909 & 0.691 \\
System 2 (+pixels) & \textbf{0.909} & \textbf{0.812} & \textbf{0.945} & \textbf{0.891} \\
System 3 (per-view)& 0.873 & 0.746 & 0.909 & 0.836 \\
\bottomrule
\end{tabular}
\caption{Development-set results ($n=55$), official evaluator. The
\textit{uniform\_random} macro figure is a single seed and is not stable; see
Section~\ref{sec:setup}.}
\label{tab:dev}
\end{table}

\begin{table}[t]
\centering\small
\begin{tabular}{lcc}
\toprule
Metric & $\Delta$ & 95\% CI \\
\midrule
view\_accuracy       & $+0.164$ & $[+0.073,+0.273]$ \\
view\_macro\_accuracy& $+0.297$ & $[+0.078,+0.461]$ \\
answer\_accuracy     & $-0.018$ & $[-0.055,+0.000]$ \\
joint\_accuracy      & $+0.146$ & $[+0.036,+0.255]$ \\
\bottomrule
\end{tabular}
\caption{Paired bootstrap, System~2 minus System~1, 4000 resamples. View
metrics improve significantly; answering does not improve.}
\label{tab:paired}
\end{table}

Table~\ref{tab:dev} gives development results. System~2 improves view
macro-accuracy from $0.143$ (constant baseline) to $0.679$, and joint accuracy
from $0.182$ to $0.855$.

Table~\ref{tab:paired} isolates the contribution of the pixels. The view
metrics improve with intervals excluding zero. Answer accuracy does
\emph{not} improve; its point estimate is negative and its interval has upper
bound exactly $0.000$.

\paragraph{The answer head is nearly saturated without vision.} System~1
reaches $0.9455$ answer accuracy having seen no image. Taken with
Table~\ref{tab:paired}, this indicates the four options are separable from
text alone, and that the remaining headroom is essentially all in view
attribution. Note that System~1's joint accuracy exactly equals its view
accuracy: answer accuracy is high enough that joint is determined by the view
head.

\paragraph{Official test results.} Table~\ref{tab:test} reports every
GoldenViewVQA run we submitted to the evaluator Space. Our best submission
reaches $0.7797$ joint accuracy with $0.8305$ view accuracy and $0.6522$ view
macro-accuracy.

\begin{table}[t]
\centering\small
\begin{tabular}{lcccc}
\toprule
Submitted run & View & Macro & Ans. & Joint \\
\midrule
constant baseline & 0.780 & 0.167 & 0.237 & 0.203 \\
System 2, GPT-5.6 Sol & 0.780 & 0.562 & \textbf{0.932} & 0.763 \\
System 2, Claude Opus 5 & \textbf{0.831} & \textbf{0.652} & 0.915 & \textbf{0.780} \\
\bottomrule
\end{tabular}
\caption{Official \texttt{test} results ($n=59$) from the GroundLM 2026
evaluator. Runs are listed in submission order.}
\label{tab:test}
\end{table}

Submitting the constant baseline to the official evaluator turned out to be
informative beyond a sanity check, because it lets us recover properties of the
withheld test labels. An always-\textsc{cam\_front}, always-\texttt{A}
submission scored $0.779661$ view accuracy, $0.237288$ answer accuracy and
$0.166667$ macro. Since $59 \times 0.779661 = 46$ and
$59 \times 0.237288 = 14$, exactly $46$ of the $59$ test records have gold view
\textsc{cam\_front} and $14$ have gold answer \texttt{A}. And because macro
accuracy averages recall over the classes present in gold, a macro of exactly
$1/6$ shows that \emph{six} of the seven view labels occur in test, not all
seven.

Two consequences follow. First, test is more front-dominated than development
($46/59 = 78.0\%$ versus $38/55 = 69.1\%$), so our systems' front-heavy
predictions were better matched to test than the development distribution
suggested. Second, one label is absent from test gold entirely; if that label
is \textsc{none\_of\_the\_above}, then our systems' inability to predict it
--- which costs $14.3$ macro points on development --- costs nothing on test,
and the macro ceiling is correspondingly higher than we assumed.

% =============================================================================
\section{Error Analysis}
\label{sec:analysis}

\begin{table}[t]
\centering\small
\begin{tabular}{lccc}
\toprule
Gold view & $n$ & Sys.\ 1 & Sys.\ 2 \\
\midrule
\textsc{cam\_front}        & 38 & 0.921 & 1.000 \\
\textsc{cam\_front\_left}  &  8 & 0.250 & 0.750 \\
\textsc{cam\_front\_right} &  2 & 0.500 & 1.000 \\
\textsc{cam\_back}         &  1 & 1.000 & 1.000 \\
\textsc{cam\_back\_left}   &  2 & 0.000 & 0.000 \\
\textsc{cam\_back\_right}  &  1 & 0.000 & 1.000 \\
\textsc{none\_of\_the\_above} & 3 & 0.000 & 0.000 \\
\bottomrule
\end{tabular}
\caption{Per-class view recall on development.}
\label{tab:perclass}
\end{table}

\paragraph{Evidence-selection failures.} Table~\ref{tab:perclass} localises
the remaining error. System~2 is perfect on four of seven classes. All residual
loss is concentrated in \textsc{cam\_back\_left} ($0/2$) and
\textsc{none\_of\_the\_above} ($0/3$), together $28.6$ macro points. Both
systems predicted \textsc{none\_of\_the\_above} zero times in every run,
despite prompts naming it as a valid and rare label. Both
\textsc{cam\_back\_left} errors went to \textsc{cam\_front\_left}: a
left-side front/rear confusion, not a left/right confusion.

\paragraph{Attention collapse.} System~1 predicted \textsc{cam\_front} on
$44/55$ development records against a gold count of $38$, and lost $6$ of $8$
\textsc{cam\_front\_left} cases to it. Adding pixels reduced this collapse
substantially, which is the main mechanism behind the macro gain.

\paragraph{Class blindness dominates, not recall.} We initially expected
detector recall to be the binding constraint. It is not. At threshold $0.15$
the front-end yields $29.3$ raw detections per frame and
$22.0$ associated objects per frame on development, against roughly $27$
annotated agents per frame in nuScenes; no frame came out empty. Instead the
binding constraint is that a seven-class agent detector cannot represent the
static infrastructure most questions turn on. Across our records, decisive
evidence was repeatedly a lowered car-park gate arm, a stop sign, a traffic
signal state, a splitter island, cones, chevron boards, a painted lane arrow or
a crosswalk --- none of which exist in the detector's vocabulary. This is a
representational ceiling, not a recall ceiling, and it explains why System~1
plateaus well below System~2.

\paragraph{Serialisation faults.} Inspecting the pixel-versus-text
disagreements our System~2 runs reported surfaced three distinct faults at the
perception-to-text interface, in descending order of severity.

\textit{(i) False negative rendered as confirmed negative.} Several cameras
were serialised as \texttt{no agents detected} while the image contained
vehicles, in one case a truck bed filling the entire frame at close range.
This is the most damaging fault available in our design, because the full
render presents an empty camera as grounds to rule that view out. We suspect
boxes whose bottom edge reaches the image border are dropped by the
ground-plane projection; clamping and flagging the object as
\textit{very near, range uncertain} would be preferable to silently emitting a
negative.

\textit{(ii) View-set misattribution.} In at least one case, objects fully
visible in \textsc{cam\_front} were attributed only to
\textsc{cam\_front\_left}, where they were clipped at the frame edge. The
fixed $2.0$\,m association radius is likely too tight at range, where
flat-ground projection error grows quadratically; a range-dependent radius is
the natural fix. This fault matters more than a false positive because view
attribution is precisely what is scored.

\textit{(iii) Phantom objects.} Cars parked in multi-storey car parks, vehicles
on the far side of a kerbed median, and a person seated on a grass embankment
were all projected into the scene as road agents. Under System~2 the pixels
arbitrate; under System~1 there is no confidence signal a text-only reader
could use to discount them.

\paragraph{Vocabulary mismatch.} The detector's seven class names do not align
with the question vocabulary. One record asks about a ``bicycle'' that the
detector labels a motorcycle; another concerns an SUV that the text calls a car
and the option calls a van. Matching on class names is therefore unreliable
even when the detection itself is correct.

\paragraph{Irreducible ambiguity.} A recurring pattern is that the decisive
object is clipped at the frame edge of one camera while fully visible in
another, so more than one view is defensible while gold names one. We expect
this to be a floor on achievable view accuracy rather than a fixable
deficiency.

\paragraph{Development versus test.} Official test performance is below
development, as anticipated, but the size of the gap differs sharply between our
two submitted runs and that difference is the more interesting result. The
Claude Opus~5 run moved from $0.679$ macro on development to $0.652$ on test, a
drop of $0.027$ that sits far inside the development interval. The GPT-5.6 Sol
run moved from $0.812$ to $0.562$, a drop of $0.250$. The run with the stronger development score
generalised markedly worse, which is what selecting on a 55-record development
set will do: the apparent development advantage was substantially
split-specific. Under the official metric the ordering of our two runs reverses
between splits.

Test is also harder in composition. It leans more on
\textit{intent\_prediction} ($20/59$ versus $13/55$), our weakest group on
development ($0.692$ view accuracy, against $0.926$ on causality). Working
against that, the recovered label statistics show test is more front-dominated
than development, which favours our systems' front-heavy predictions. Our
submitted run predicted \textsc{cam\_front} on $47/59$ records against a gold
count of $46$, so the front-heaviness we had flagged as a risk turned out to be
close to calibrated rather than a systematic error.

% =============================================================================
\section{Conclusion}
\label{sec:conclusion}

Serialising a surround-view scene into text with explicit per-object view-sets
turns view-level evidence identification into a grounding-and-lookup problem
and takes view macro-accuracy from $0.143$ to $0.382$ with no visual input.
Supplying the raw images on top of that serialisation nearly doubles it again
to $0.679$, and a paired bootstrap confirms the gain is in the view head only.

The most useful thing we can report is negative. Answer accuracy of $0.9455$
without any image indicates that the answer options are largely separable from
text alone, and that joint accuracy therefore behaves as a view metric. Future
work on this benchmark should concentrate on view attribution, and we would
encourage the organizers to publish a question-and-options-only baseline so
that the answer-side leak is quantified for all participants.

A second lesson is methodological. Our two models' rankings reverse between
splits: the run that scored $0.812$ macro on development fell to $0.562$ on
test, while the weaker development run ($0.679$) held at $0.652$. On a
55-record development set, selecting a submission by development score is
close to selecting noise, and we would have submitted the wrong run had we
trusted it.

For our own system, the clearest next steps follow from
Section~\ref{sec:analysis}: never serialise an unverified negative, make the
association radius range-dependent, and add an explicit abstention rule so that
\textsc{none\_of\_the\_above} becomes reachable rather than structurally
impossible.

% =============================================================================
\section*{Limitations}

Our evaluation rests on 55 labelled records. Bootstrap intervals are
correspondingly wide --- $0.306$ for System~2 macro accuracy --- and two view
classes contain a single example each, so per-class recalls for those classes
are one-bit measurements. We showed that changing only a random seed moves
macro accuracy over a $0.208$ range; readers should not treat differences
smaller than that as meaningful.

The rig geometry we hard-code is nuScenes-specific and would not transfer to
another sensor platform without re-derivation. The perception front-end is
monocular and assumes flat ground, which is the direct cause of the
phantom-object and misattribution faults documented above.

Our detector was fine-tuned on the nuScenes \texttt{v1.0-mini} split, whose ten
scenes include \texttt{scene-0061} and \texttt{scene-0103}. Both appear in
GoldenViewVQA, contributing one development record
(\texttt{sfall\_0103\_causality}) and one test record
(\texttt{sfall\_0061\_intent\_prediction}).

Our systems consume exactly the six images of the annotated keyframe and
nothing else. Because each record identifies a nuScenes sample within a
continuous scene, temporally adjacent frames are available, and a
similarity-based retrieval step --- gathering nearby frames whose scene content
is close to the queried keyframe and serialising them jointly --- would give
substantially richer context than a single instant. This would particularly
help the \textit{intent\_prediction} group, our weakest at $0.692$ view
accuracy, where a pedestrian's direction of travel or a vehicle's manoeuvre is
far more legible over a short window than in one frame. It would also supply
the motion cues our single-frame design discards by construction.

Both submitted systems depend on a cloud-hosted multimodal model. Beyond cost
and reproducibility concerns --- a hosted model may be updated or withdrawn,
so our exact numbers may not be reproducible later --- the latency is
disqualifying for the application domain this benchmark is drawn from. Each
System~2 record requires transmitting six full-resolution images and waiting on
a remote inference call, which is orders of magnitude away from the budget of
an on-vehicle perception loop. A safety-relevant evidence-attribution component
would have to run locally. Investigating smaller open-weight vision-language
models, and distilling the behaviour of the larger model into one, is the more
practical direction, and it is the direction our prior work took for its
reasoning module \citep{gaafar2025geodrive}. We expect a meaningful accuracy
cost, and quantifying that trade-off is exactly what makes the experiment
worth running.

\section*{Ethics Statement}

Images were obtained from the official nuScenes distribution under its Terms
of Use. No images are redistributed with our code or in this paper, consistent
with the shared task's own policy. We use no ground-truth nuScenes annotations
at inference, and no personal data beyond what the public nuScenes release
already contains; pedestrians appear in the source imagery as distributed by
nuScenes, and we perform no identification of individuals.

All external resources are disclosed in Section~\ref{sec:data} and
Section~\ref{sec:system}: the nuScenes dataset, a YOLOv8-m checkpoint
fine-tuned on nuScenes from our prior work, and a commercial multimodal
language model accessed through its provider. No synthetic or generated
training data was used, and no model was trained for this submission.

% =============================================================================
\bibliography{refs}

\end{document}
